\pagestyle{plain}

\begin{center}
    \textbf{\textit{Streszczenie}}
\end{center}

Niniejsza praca magisterska poświęcona jest analizie porównawczej wybranych nierelacyjnych baz danych typu klucz-wartość (Redis, Memcached, LevelDB, RocksDB, Couchbase, Tarantool, FoundationDB oraz BerkeleyDB) w środowisku zwirtualizowanym WSL 2 oraz natywnym systemie Linux. Głównym celem pracy było zbadanie wpływu architektury silnika bazy danych oraz modelu wdrożenia na przepustowość, opóźnienia oraz efektywność zasobową w zróżnicowanych scenariuszach obciążenia (operacje CRUD, obciążenia mieszane, operacje na dokumentach JSON oraz strukturach listowych). W ramach badań zaprojektowano i zaimplementowano aplikację benchmarkującą w języku Python, która pozwoliła na symulację obciążenia wieloprocesowego. Wyniki pomiarów poddano szczegółowej analizie statystycznej, uwzględniając odchylenia standardowe oraz opóźnienia ogonowe (centyle p95 i p99). Pozwoliło to na sformułowanie praktycznych rekomendacji wdrożeniowych dotyczących wyboru odpowiedniego silnika bazodanowego w zależności od specyfiki realizowanego projektu.

\par\vspace{0.5cm}
\noindent
\textbf{Słowa kluczowe:} \textit{bazy danych NoSQL, systemy klucz-wartość, benchmark, WSL 2, Linux, wydajność baz danych, opóźnienia ogonowe}

\vspace{1.5cm}

\begin{center}
    \textbf{\textit{Abstract}}
\end{center}

This master's thesis presents a comparative performance analysis of selected non-relational key-value databases (Redis, Memcached, LevelDB, RocksDB, Couchbase, Tarantool, FoundationDB, and BerkeleyDB) in the WSL 2 virtualized environment and the native Linux operating system. The primary objective of the study was to evaluate the impact of database engine architecture and deployment model on throughput, latency, and resource efficiency under diverse workload scenarios (CRUD operations, mixed workloads, JSON document manipulations, and list-based operations). A dedicated benchmarking application was designed and implemented in Python to simulate multi-process client workloads. The experimental results were subjected to detailed statistical analysis, focusing on standard deviations and tail latencies (p95 and p99 centiles). This enabled the formulation of practical architectural recommendations for selecting the optimal database engine based on specific system requirements.

\par\vspace{0.5cm}
\noindent
\textbf{Keywords:} \textit{NoSQL databases, key-value systems, benchmark, WSL 2, Linux, database performance, tail latency}
